\documentclass[11pt,a4paper]{article}
\usepackage[T1]{fontenc}\usepackage[utf8]{inputenc}\usepackage{lmodern}
\usepackage[a4paper,margin=2.2cm]{geometry}
\usepackage{amsmath,amssymb}\usepackage{enumitem}\usepackage{booktabs,longtable,array}
\usepackage[hidelinks]{hyperref}
\setlength{\parindent}{0pt}\setlength{\parskip}{0.4em}
\title{\textbf{OP3 / B2-Q1 audit: does the recorded $n_s$ conversion step
require the split-branch window?}\\\large Recon v1 --- proposal-only}
\author{Per reviewer directive after acceptance of B2 rev.~1.1}
\date{12 June 2026}
\begin{document}\maketitle

\section*{0. Scope and firewall}
B2-Q1 asked whether the recorded critical-fluctuation route to
$n_s\sim0.96$--$0.97$ presupposes the split-branch window (potential-imported
gap) or is formulated where the B2 sublemma is silent. This audit answers it
by dissecting the recorded construction (sec:primordial\_estimates) stage by
stage against the recorded interface discipline. Proposal-only; no preprint
edits; conditional on the Level-B kernel and on the recorded
Planck-threshold matching \emph{estimate} (the record's own wording).

\section{The recorded conversion architecture, dissected}
The recorded route has three stages with different homes:
\begin{description}[leftmargin=2.0em,itemsep=2pt]
\item[Stage 1 (pre-emergent).] The critical correlator: $O(4)$ criticality
at the upper critical dimension, $G_4(p)$ with logarithmic corrections
$G(r)\sim r^{-2}[\ln(r/\xi_0)]^{-\hat p}$. This is Euclidean critical
statistics of the order parameter --- no emergent time, no metric, no
energy. The B2 sublemma is silent here by construction.
\item[Stage 2 (hand-off at onset).] The fluctuation is read as a local
transition-time imprint, $\zeta(\mathbf x)\sim\delta u_*/u_0$: critical
statistics become \emph{initial data} for the ordered branch at onset. The
recorded thermal-completion discipline fixes the reading: onset quantities
are ``defined after emergence/matching within the FLRW reduction --- not an
energy statement at the pre-emergent level.''
\item[Stage 3 (post-onset).] The conversion proper,
$k^2\Phi_k\sim G_N\delta\rho_k$, and the gauge-disciplined freeze-out. Both
use post-onset structure: an FLRW reduction, a matched $G_N$, an emergent
gauge sector. By the record's own discipline these live strictly on the
ordered side.
\end{description}

\section{Where the ordered EFT is actually evaluated}
The only place the ordered EFT is invoked is Stage 3, at or after branch
onset. Under the recorded Planck-threshold matching estimate the onset
curvature is $K_{\rm onset}\sim H_{\rm onset}^2$ with
$|H_{\rm onset}|^{-1}$ of order $M_{\rm Pl}^{-1}$, while the gap at the
matched point is condensate-scale: the hand-off is \emph{marginal},
$K_{\rm onset}\sim m_\sigma^2$, on \emph{both} F7 branches, and the
description improves monotonically afterwards as $K$ falls and $u$ settles
to its track. The small-$u$ region \emph{before} onset --- the only region
where the merged and split branches differ --- is never used by the
recorded route: it is covered by Stage 1, which is pre-emergent.

\section{Audit verdict}
\begin{enumerate}[leftmargin=1.9em,itemsep=2pt,label=\textbf{A\arabic*.}]
\item \textbf{B2-Q1 answered: the recorded $n_s$ route does not presuppose
the split window.} The construction is hand-off-at-onset by its own
recorded discipline: pre-emergent critical statistics (where B2 is silent)
$\to$ initial data at a marginal onset boundary $\to$ post-onset conversion
inside the validity domain. It is therefore compatible with \emph{both} F7
branches; the merged (gradient-native) branch remains compatible with
(is not excluded by) the
recorded observable estimate.
\item \textbf{Named residue: hand-off marginality.} Because
$K_{\rm onset}\sim m_\sigma^2$ at the matching estimate, the Stage-2/3
hand-off carries $O(1)$ conversion uncertainty. This does \emph{not}
touch the tilt's structure --- the logarithmic form
$1-n_s\sim\hat p/\ln(L/\xi_0)$ is fixed in Stage 1 --- but it does sit on
top of the amplitude bookkeeping ($\delta u_*/u_0\sim5\times10^{-5}$),
which already carries an order-of-magnitude conversion/coarse-graining
factor in the record. A one-line acknowledgement of this $O(1)$ hand-off
factor in the recorded Status paragraph is a D-candidate --- only after
review, and only if judged reader-useful.
\item \textbf{Consequence for the discriminator.} The merged/split
dichotomy is \emph{not} resolved by the recorded $n_s$ route: the
discriminator remains a transition-side structural question, exactly as
the late-time-invisibility note of B2 anticipated. No recorded observable
currently distinguishes the F7 branches.
\item \textbf{Honesty clause.} All statements are conditional on the
Level-B kernel, on the F7 assumption set, and on the Planck-threshold
matching being what the record calls it --- an estimate. If onset occurs
parametrically below the threshold ($K_{\rm onset}\ll m_\sigma^2$), the
hand-off ceases to be marginal and A2 weakens to a smaller correction;
the A1 conclusion is unaffected in either case.
\end{enumerate}
\end{document}
